% !TEX program = xelatex
% 阶段工作入口：目标、实现、实验、问题与计划。
\documentclass[aspectratio=169,10pt,fontset=none]{ctexbeamer}
\usetheme{BITUniversal}
% 所有入口共用。复制项目后，先改这里的标题、姓名、机构和日期。
\title[组会汇报]{组会汇报}
\subtitle{研究主题 / 论文中文标题}
\setenglishtitle{Research Topic / Original Paper Title}
\setpresenter{汇报人姓名}
\setaffiliation{北京理工大学}
\setmeeting{课题组组会}
\date{\number\year.\ifnum\month<10 0\fi\number\month.\ifnum\day<10 0\fi\number\day}
% \date{2026.11.09}  % 如需固定日期，取消本行注释。
\setpaper{Original Paper Title}{Author A, Author B, et al.}{Conference / Journal, Year}

% 默认：DT 参考中的轻量白底页眉；改成 band 使用 PWBC 绿色渐变页眉。
\setheaderstyle{light}
% \bitsectionpagesfalse  % 短汇报可关闭自动章节过渡页。

\begin{document}
\BITtitlepage
\BITcontents
\bitsection{目标与当前问题}{本阶段任务与上次组会的后续工作}
\begin{frame}{阶段目标与完成情况}
  \small\renewcommand{\arraystretch}{1.8}
  \begin{tabularx}{\textwidth}{@{}XXX@{}}
    \toprule\bitstrong{阶段任务} & \bitstrong{当前状态} & \bitstrong{结果证据}\\\midrule
    \bitplaceholder{任务 A} & \bitplaceholder{已完成 / 进行中} & \bitplaceholder{结果或产物}\\
    \bitplaceholder{任务 B} & \bitplaceholder{已完成 / 进行中} & \bitplaceholder{结果或产物}\\\bottomrule
  \end{tabularx}
\end{frame}
\bitsection{方案与实现}{本阶段的具体改动与技术依据}
\begin{frame}{方案与关键改动}
  \bitimageplaceholder[.46\textheight]{插入当前方案、系统结构或实现截图}
  \bitcallout{本阶段改动：\bitplaceholder{改了什么，为什么这样改。}}
\end{frame}
\bitsection{实验与问题分析}{设置、结果、异常现象与排查证据}
\begin{frame}{实验设置与结果}
  \bitexperiment{\bitplaceholder{验证当前改动的效果。}}{\bitplaceholder{基线、控制变量与统计方式。}}
  \bitimageplaceholder[.43\textheight]{插入结果曲线、定量表格或行为对比}
  \bitcallout{观察：\bitplaceholder{陈述有证据支持的结果。}}
\end{frame}
\begin{frame}{当前问题与排查结果}
  \begin{description}
    \item[现象] \bitplaceholder{异常在哪种条件下发生。}
    \item[已有证据] \bitplaceholder{哪些实验缩小了原因范围。}
    \item[待验证假设] \bitplaceholder{明确标注尚未确认的解释。}
  \end{description}
\end{frame}
\bitsection{计划与讨论}{下一步任务、完成标准与需要讨论的问题}
\begin{frame}{后续计划}
  \small\renewcommand{\arraystretch}{1.8}
  \begin{tabularx}{\textwidth}{@{}p{22mm}XX@{}}
    \toprule\bitstrong{时间} & \bitstrong{任务} & \bitstrong{完成标准}\\\midrule
    近期 & \bitplaceholder{优先解决的问题} & \bitplaceholder{可检查的结果}\\
    下一阶段 & \bitplaceholder{后续实验或实现} & \bitplaceholder{预期产出}\\\bottomrule
  \end{tabularx}
  \bitcallout{需要讨论：\bitplaceholder{具体的技术选择或实验安排。}}
\end{frame}
\BITthankyou
\end{document}
